Example: 1) $S_4$
No. of cycles of $2$ length in $S_4$: $\frac{4!}{(1! 2^1)(2! 1^2)}$
$= \frac{4 \times 3 \times 2}{2 \times 2}$
$= 6$

2) $S_5$
No. of cycle of length $2$: $\frac{5!}{(3! 1^3)(1! 2^1)} = \frac{5 \times 4 \times 3 \times 2 \times 1}{3 \times 2 \times 1 \times 2 \times 1}$
$= \frac{20}{2} = 10$

$\frac{5!}{(1! 2^1)(1! 3^1)} = \frac{5 \times 3 \times 2 \times 4 \times 1}{2 \times 3}$
$= 20$

# Ring: A system of $(R, +, \cdot)$ consisting of a non-empty set $R$ and two internal binary operation '+' $\&$ '$\cdot$' in $R$ is called a ring if
1) $(R, +)$ is an abelian group
2) $(R, \cdot)$ is a semi group.
3) $a(b + c) = a \cdot b + a \cdot c \quad \forall a, b, c \in R$
$(a + b) \cdot c = a \cdot c + b \cdot c$

If
- wit $(R, \cdot)$ has identity $\Rightarrow$ Ring with unity
- If $(R, \cdot)$ has is abelian $\Rightarrow$ Ring Commutative ring
- Identity of $(R, +) \Rightarrow$ zero element of ring.

Note: 1) '+' is usually called addition and
2) '$\cdot$' is usually called multiplication.
3) Identity w.r.t. '+' is called as the zero element of ring and is denoted as $0$.
